% id: 38-24
% book: 베이직쎈 미적분1 · 02 함수의 연속
% section: 개념 18 사잇값 정리
% figure: none
% answer: 
% answer_source: 
% variant_level: 0 (원본 전사)
% 이 파일은 build.mjs 가 items.json 에서 만든다. 고칠 때는 items.json 을 고친다.
\smash{\raisebox{1.5mm}{\dmkichul{베이직쎈 미적분1 38쪽 24번}}}
다음은 방정식 $x^3+x+3=0$이 열린구간 $(-2{,}\mkern3mu\mathopen{}0)$에서 적어도 하나의 실근을 가짐을 증명한 것이다. \blank\ 안에 알맞은 것을 써넣으시오.\exprbox{\begin{minipage}{0.96\linewidth}{\small\bfseries 증명}\par\smallskip\raggedright $f(x)=x^3+x+3$이라 하면 함수 $f(x)$는 닫힌구간 $[-2{,}\mkern3mu\mathopen{}0]$에서 연속이다.\par 또 $f(-2)=\blank$, $f(0)=\blank$에서\par\smallskip\dispeq{$f(-2)f(0)\mkern3mu \blank\mkern3mu 0$}\par\smallskip 이므로 사잇값 정리에 의하여 $f(c)=\blank$인 $c$가 열린구간 $(-2{,}\mkern3mu\mathopen{}0)$에 적어도 하나 존재한다.\par 따라서 방정식 $x^3+x+3=0$은 열린구간 $(-2{,}\mkern3mu\mathopen{}0)$에서 적어도 하나의 실근을 갖는다.\end{minipage}}
